\documentclass[12pt,oneside]{book}

% ==================================================================
% METADATOS
% ==================================================================
\newcommand{\universidad}{Universidad de Buenos Aires (UBA)}
\newcommand{\facultad}{Facultad de Derecho}
\newcommand{\materia}{Derecho Constitucional Profundizado}
\newcommand{\titulo}{Apuntes de Derecho Constitucional Profundizado}
\newcommand{\subtitulo}{Libertad religiosa \textperiodcentered{} Categorías sospechosas \textperiodcentered{} Defensa del consumidor \textperiodcentered{} \textit{Solve et repete} y lealtad comercial}
\newcommand{\autor}{Isaac Edgar Camacho Ocampo}
\newcommand{\anio}{2026}

% ==================================================================
% PAQUETES
% ==================================================================
\usepackage[utf8]{inputenc}
\usepackage[T1]{fontenc}
\usepackage[spanish]{babel}

\usepackage[
    top=2.5cm,
    bottom=2.5cm,
    left=3cm,
    right=2.5cm,
    headheight=15pt
]{geometry}

\usepackage{amsmath}
\usepackage{amssymb}
\usepackage{mathtools}

\usepackage{graphicx}
\usepackage{float}
\usepackage{caption}
\usepackage{subcaption}

\usepackage{booktabs}
\usepackage{array}
\usepackage{longtable}
\usepackage{tabularx}

\usepackage{enumitem}

\usepackage{xcolor}
\usepackage[most]{tcolorbox}

\usepackage{fancyhdr}
\usepackage{ragged2e}

% ==================================================================
% CONFIGURACIÓN
% ==================================================================
\graphicspath{
    {./img/}
    {./graficos/}
    {./figuras/}
}

\setcounter{tocdepth}{2}

\setlength{\parindent}{0pt}
\setlength{\parskip}{0.5em}

\setlist[itemize]{
    topsep=0.3em,
    itemsep=0.2em
}

\setlist[enumerate]{
    topsep=0.3em,
    itemsep=0.2em
}

\clubpenalty=10000
\widowpenalty=10000

\pagestyle{fancy}
\fancyhf{}
\fancyhead[L]{\nouppercase{\leftmark}}
\fancyhead[R]{\materia}
\fancyfoot[C]{\thepage}
\renewcommand{\headrulewidth}{0.4pt}
\renewcommand{\footrulewidth}{0pt}

\fancypagestyle{plain}{
    \fancyhf{}
    \fancyfoot[C]{\thepage}
    \renewcommand{\headrulewidth}{0pt}
}

% ==================================================================
% COMANDOS
% ==================================================================
\newcommand{\abs}[1]{\left|#1\right|}
\newcommand{\norm}[1]{\left\lVert#1\right\rVert}
\newcommand{\vect}[1]{\mathbf{#1}}
\newcommand{\deriv}[2]{\frac{d#1}{d#2}}
\newcommand{\pderiv}[2]{\frac{\partial#1}{\partial#2}}

% ==================================================================
% ENTORNOS / CAJAS
% ==================================================================
\newtcolorbox{definition}[1]{
    enhanced,
    breakable,
    title={Definición: #1},
    fonttitle=\bfseries,
    colback=blue!3,
    colframe=blue!50!black,
    boxrule=0.8pt,
    arc=2mm
}

\newtcolorbox{important}{
    enhanced,
    breakable,
    title={Importante},
    fonttitle=\bfseries,
    colback=yellow!8,
    colframe=orange!70!black,
    boxrule=0.8pt,
    arc=2mm
}

\newtcolorbox{example}[1]{
    enhanced,
    breakable,
    title={Ejemplo: #1},
    fonttitle=\bfseries,
    colback=green!4,
    colframe=green!40!black,
    boxrule=0.8pt,
    arc=2mm
}

\newtcolorbox{formula}[1]{
    enhanced,
    breakable,
    title={Fórmula / Regla: #1},
    fonttitle=\bfseries,
    colback=purple!4,
    colframe=purple!50!black,
    boxrule=0.8pt,
    arc=2mm
}

\newtcolorbox{exercise}[1]{
    enhanced,
    breakable,
    title={Ejercicio: #1},
    fonttitle=\bfseries,
    colback=gray!5,
    colframe=gray!60!black,
    boxrule=0.8pt,
    arc=2mm
}

\newtcolorbox{note}{
    enhanced,
    breakable,
    title={Nota},
    fonttitle=\bfseries,
    colback=white,
    colframe=gray!50,
    boxrule=0.6pt,
    arc=2mm
}

% ==================================================================
% CONFIGURACIÓN FINAL
% ==================================================================
\usepackage[
    colorlinks=true,
    linkcolor=blue!50!black,
    citecolor=green!40!black,
    urlcolor=blue!60!black,
    pdftitle={\titulo},
    pdfauthor={\autor},
    pdfsubject={Apuntes universitarios},
    bookmarks=true
]{hyperref}

% ==================================================================
% DOCUMENTO
% ==================================================================
\begin{document}

% ------------------------------------------------------------------
% PORTADA
% ------------------------------------------------------------------
\begin{titlepage}

\thispagestyle{empty}

\begin{center}

\vspace*{1cm}

{\large \textsc{\universidad}}

\vspace{0.2cm}

{\large \textsc{\facultad}}

\vspace{3cm}

{\Huge \textbf{\materia}}

\vspace{1cm}

{\LARGE \titulo}

\vspace{0.4cm}

{\normalsize \subtitulo}

\vspace{0.8cm}

\textit{Comprender es el primer paso para convertir el estudio en conocimiento.}

\vspace{1.2cm}

\rule{0.70\textwidth}{0.5pt}

\vspace{0.5cm}

{\large Apuntes de estudio}

\vspace{0.5cm}

\rule{0.70\textwidth}{0.5pt}

\vfill

{\large \autor}

\vspace{0.3cm}

{\large \anio}

\end{center}

\vfill

\begin{flushleft}
\textit{Este es un modesto aporte a los estudiantes de ``Derecho Constitucional Profundizado'' no pretende sustituir el material oficial de la materia.}
\end{flushleft}

\end{titlepage}

% ------------------------------------------------------------------
% ÍNDICE
% ------------------------------------------------------------------
\tableofcontents

% ------------------------------------------------------------------
% INTRODUCCIÓN
% ------------------------------------------------------------------
\chapter*{Introducción}
\addcontentsline{toc}{chapter}{Introducción}
\markboth{Introducción}{Introducción}

Estos apuntes recorren cuatro ejes que comparten una misma idea: \textbf{el poder, ya sea del Estado o de las empresas, tiene límites}. Quien lo ejerce debe tratar a todos con igualdad, razonabilidad y garantías.

\begin{enumerate}
    \item La relación entre el Estado y la Iglesia y la \textbf{libertad religiosa}, a partir de un fallo de la Corte contra la provincia de Salta.
    \item Las \textbf{categorías sospechosas} y la inversión de la carga de la prueba, con ejemplos de edad, nacionalidad y género.
    \item El \textbf{poder sancionatorio} en la defensa del consumidor y las garantías del debido proceso.
    \item El \textit{solve et repete} y la \textbf{lealtad comercial}.
\end{enumerate}

Cada capítulo concluye con un cuadro Cornell para el repaso mediante recuperación activa: conviene intentar responder cada pregunta de la columna izquierda antes de leer la columna derecha.

% ==================================================================
% CAPÍTULO 1
% ==================================================================
\chapter[Libertad religiosa y Estado--Iglesia]{Libertad religiosa y relación Estado--Iglesia}

\section{Herencia colonial y supresión del requisito religioso (1994)}

\subsection{El culto predominante y la praxis jurídica}

Existe una ligazón entre el orden jurídico, los actores principales de la praxis jurídica y el culto religioso predominante. Se trata de una herencia de la época colonial hispánica.

\subsection{El requisito del culto católico y su supresión}

La Constitución exigía, para acceder al cargo de presidente, ser practicante del culto católico apostólico romano. Este requisito fue uno de los puntos principales de la temática que se modificó: \textbf{se suprimió directamente en el año 1994}. Con ello se afianzó la idea de que el culto no debía influir sobre la gestión política ni sobre la gestión institucional del Estado.

% TODO: contenido incompleto o ambiguo en las notas originales (la frase sobre el cargo exigido aparece reconstruida entre corchetes en la transcripción)

Las modificaciones introducidas al suprimir este punto no fueron tan sustanciales. Otras reformas provocaron mayor resistencia de la Iglesia. Sin embargo, el requisito dejó de ser habilitante, y la Convención Constituyente manifestó un fuerte arraigo del sostenimiento total de la \textbf{igualdad religiosa} y del \textbf{trato igualitario} de las religiones oficialmente reconocidas, es decir, aquellas inscriptas en el registro de cultos de la Secretaría.

\begin{example}{El caso de Menem}
En clase se comentó que Menem, hijo de inmigrantes sirios, tenía una tradición musulmana dentro de su familia, y que públicamente, desde mucho antes de su candidatura presidencial, manifestaba ser católico. Las motivaciones o la índole personal en estos casos son muy complejas, y legislar sobre la pertenencia religiosa de los funcionarios es una materia bastante inasible.
\end{example}

\begin{important}
En 1994 se suprimió el requisito religioso para acceder al cargo: la pertenencia al culto católico dejó de ser un requisito habilitante, y la Convención Constituyente reafirmó la igualdad religiosa y el trato igualitario de los cultos oficialmente reconocidos.
\end{important}

\section{El caso de Salta: la demanda y la norma provincial}

\subsection{La demanda de una asociación contra la provincia}

El fallo analizado reseña la posición actual de la Corte. Parte de una demanda de una asociación contra la provincia de Salta por la educación religiosa en los colegios.

La Constitución de la provincia establecía que todos los alumnos de la escuela primaria tenían derecho a recibir, en el seno de la educación elemental estatal, la educación de la religión que profesen. No reconocía expresamente que debiera o pudiera impartirse catecismo, es decir, religión católica, en la educación primaria. Lo que establecía era el acceso a la educación del culto que cada uno profesase.

\subsection{Qué se atacaba en la demanda}

La demanda solicitaba la declaración de inconstitucionalidad, pero cuestionaba \textbf{cómo se aplicaba en la praxis provincial} esa norma. No se atacaba el artículo en sí mismo como inconstitucional, sino su aplicación. La Corte falló casi por unanimidad, con una disidencia parcial del juez Rosatti.

\begin{important}
La demanda no atacó el texto de la norma provincial, sino su aplicación en la práctica.
\end{important}

\section{Doctrina de la Corte: neutralidad, artículos 14 y 2}

\subsection{El cambio de 1994 y el sostenimiento del culto}

La Corte destaca el cambio del año 1994 como un cambio significativo en aras de asegurar un trato equitativo y una dimensión objetiva de lo que es el sostenimiento del culto. Por pertenencia cultural y raíces históricas, existía un apoyo institucional y fundamentalmente económico al culto católico.

\begin{quote}
\textbf{Art. 2, Constitución Nacional (texto original no reformado).} El Gobierno federal sostiene el culto católico apostólico romano.
\end{quote}

\subsection{El estándar constitucional: igualdad y neutralidad}

Como estándar constitucional en torno a la libertad religiosa, sostenido por la Corte en este fallo y en fallos anteriores, el sostenimiento económico y el reconocimiento institucional derivados de las raíces históricas quedan limitados a ese posicionamiento. No suponen un trato inequitativo ni una deferencia, y mucho menos un apoyo efectivo y operativo al culto católico.

El Estado solamente cumple con un sostenimiento económico. Frente a todos los cultos reconocidos en el país debe practicar una política de igualdad, de no discriminación y de no preferencia.

\begin{definition}{Principio de neutralidad}
Política de igualdad, de no discriminación y de no preferencia del Estado frente a todos los cultos reconocidos. Es un concepto comúnmente conocido en el sistema de la mayoría de los países europeos, y parte de una sentencia del tribunal alemán que luego fue tomado por la mayoría de los países.
\end{definition}

\subsection{Compatibilizar el artículo 14 con lo que queda del artículo 2}

Ya desde antes de este fallo, la Corte sostenía que la libertad religiosa asegurada por el Estado en el artículo 14 debía \textbf{compatibilizarse} con lo que queda del artículo 2. El artículo 2 quedaba acotado a un apoyo institucional y económico, y de ninguna forma debía entenderse como una preferencia activa, en las políticas públicas, a la difusión y al credo del culto católico.

\begin{quote}
\textbf{Art. 14, Constitución Nacional.} Todos los habitantes de la Nación gozan de los siguientes derechos conforme a las leyes que reglamenten su ejercicio; a saber: de trabajar y ejercer toda industria lícita; de navegar y comerciar; de peticionar a las autoridades; de entrar, permanecer, transitar y salir del territorio argentino; de publicar sus ideas por la prensa sin censura previa; de usar y disponer de su propiedad; de asociarse con fines útiles; \textbf{de profesar libremente su culto}; de enseñar y aprender.
\end{quote}

\begin{important}
El artículo 2 se limita a un sostén institucional y económico; el artículo 14 (libertad de culto) exige igualdad, no discriminación y neutralidad del Estado frente a todos los cultos.
\end{important}

\section{Hechos probados y decisión de la Corte}

\subsection{Lo que demostró la asociación demandante}

La demanda no cuestionaba la redacción del artículo de la Constitución provincial, que hablaba del derecho a recibir la religión que cada uno profese. En la práctica, la asociación demandante demostró lo siguiente:

\begin{itemize}
    \item La materia ``religión'' se daba en todos los cursos de la educación primaria.
    \item Para cumplir el mandato provincial de enseñar la religión que cada uno profese, las familias de los estudiantes menores de edad debían llenar un formulario indicando qué religión profesaban. Esto implicaba una obligación de exteriorizar un elemento de convicción íntima que puede entenderse comprendido en el \textbf{principio de reserva}. A modo de comparación, constituciones más nuevas, como la de Tierra del Fuego, prohíben los cuestionarios o indagaciones públicas donde deba exteriorizarse la religión.
    \item Todos los profesores estaban capacitados solo para explicar la doctrina de la Iglesia, es decir, un catecismo básico de la religión católica. En los hechos solo se daba catecismo católico, porque era lo único para lo que estaban capacitados.
    \item Existía una suerte de objeción de conciencia o religiosa: los alumnos no católicos podían no estar en la clase presencial. Pero no tenían ninguna actividad alternativa durante las horas de catecismo, que la gran mayoría cursaba. Quienes no asistían deambulaban por el colegio o tomaban un café en el bar o restaurante del establecimiento, si lo había. Eran horas perdidas, y el Ministerio de Educación no tenía ninguna política de ofrecerles una alternativa educativa, ya fuera de otras religiones o, al menos, de alguna materia.
    \item No se impartían clases ni cursos de otras religiones, ni en el colegio ni mediante expertos de otros cultos.
\end{itemize}

\begin{quote}
\textbf{Art. 19, Constitución Nacional (principio de reserva; norma aludida sin cita expresa del artículo).} Las acciones privadas de los hombres que de ningún modo ofendan al orden y a la moral pública, ni perjudiquen a un tercero, están sólo reservadas a Dios, y exentas de la autoridad de los magistrados. Ningún habitante de la Nación será obligado a hacer lo que no manda la ley, ni privado de lo que ella no prohíbe.
\end{quote}

El panorama probado en el expediente era, de hecho, casi un monopolio de clases de religión católica, sin opciones para quienes profesaran otras religiones, con obligación de exteriorizar el culto y con horas perdidas sin alternativas educativas reales a ningún nivel.

\begin{table}[H]
\centering
\caption{Lo que decía la norma frente a lo que ocurría en la práctica}
\small
\begin{tabularx}{\textwidth}{
    >{\raggedright\arraybackslash}p{0.19\textwidth}
    >{\raggedright\arraybackslash}X
    >{\raggedright\arraybackslash}X
}
\toprule
\textbf{Aspecto} & \textbf{Letra de la norma provincial} & \textbf{Práctica probada en el expediente} \\
\midrule
\textbf{Contenido de la enseñanza} & Derecho a recibir la religión que cada uno profese (general y ecuménico). & Solo se enseñaba la doctrina católica (catecismo básico). \\
\textbf{Preparación docente} & No se exigía una religión determinada. & Los profesores estaban capacitados únicamente para enseñar religión católica. \\
\textbf{Identificación de la familia} & No se hablaba de declarar la religión. & Formulario anual en el que las familias debían declarar su culto. \\
\textbf{Alumnos no católicos} & Derecho a su propia educación religiosa. & Podían no asistir, pero sin actividad sustitutiva: horas perdidas. \\
\textbf{Otras religiones} & Posibilidad abstracta de recibirlas. & No se impartían clases ni cursos de otras religiones. \\
\bottomrule
\end{tabularx}
\end{table}

\subsection{El razonamiento de la Corte}

La Corte sostiene que, aun cuando no hay a priori una inconstitucionalidad en el propio texto, la práctica sí la genera. El texto podría quedar dentro de una zona de reserva de la provincia, como un modo de disponer cómo abordar el tema religioso en las escuelas, y su redacción no indicaba favoritismo ni preferencia de un culto sobre otro. Pero las pruebas eran contundentes: las clases de religión equivalían en realidad a clases del culto católico apostólico, y no había ninguna alternativa para quienes quisieran estar por fuera de ellas.

La Corte señala que \textbf{la práctica es inconstitucional y viola el principio de neutralidad} en un aspecto clave, porque la educación es formativa en edades tempranas. Ese ecosistema educativo, en el que solo se daban contenidos católicos, exteriorizaba una conducta que atentaba contra la igualdad religiosa, era inequitativa y ponía a esta religión en un pie de desigualdad respecto de otras.

La Corte reitera su doctrina de los últimos años, mantenida en numerosos expedientes. Según esta doctrina, también por el contenido religioso puede haber una conducta de discriminación, y toda práctica que haga prevalecer una religión sobre otra entra en un marco de sospecha de discriminación. Ello no se corresponde con el estándar del artículo 14, en el cual el sostenimiento queda relegado a funciones materiales de provisión, de asistencia y de sostenimiento estrictamente económico.

Por lo tanto, aun cuando existe una zona de reserva para regular un aspecto de su competencia originaria, como la educación básica, la provincia no puede generar en la práctica una política que viole el trato no discriminatorio y el pie de igualdad entre todas las religiones. La Corte declaró inválida la forma en que se ponía en práctica el sistema, sin perjuicio de que la provincia pueda establecer, dentro de un marco razonable, las características que pretenda para la educación primaria, siempre que salvaguarde el tratamiento igualitario para todas las religiones.

\begin{important}
Una norma sin defectos aparentes en su texto puede ser inválida por su aplicación. La provincia conserva su competencia sobre la educación básica, pero debe garantizar un tratamiento igualitario a todas las religiones.
\end{important}

\section{Estatus jurídico de la Iglesia Católica}

\subsection{Persona jurídica pública y concordato}

En el Código Civil y Comercial vigente, la Iglesia Católica es una \textbf{persona jurídica de derecho público}. Es la religión mayoritaria, aunque resulta difícil censar la práctica constante de una religión. Tiene un \textbf{concordato} con estatus de tratado internacional, en el que rige un principio de cooperación y un principio de respeto a las investiduras eclesiásticas, que son designadas desde el Vaticano.

\begin{quote}
\textbf{Art. 146, Código Civil y Comercial de la Nación (norma aludida; se transcribe el inciso pertinente).} Personas jurídicas públicas. Son personas jurídicas públicas: [\ldots] c) la Iglesia Católica.
\end{quote}

\subsection{Bienes afectados al culto e inembargabilidad}

La Iglesia tiene un derecho sobre sus bienes de propiedad, que ejerce. En la medida en que esos bienes están dirigidos a impartir el culto católico, es decir, afectados a la liturgia del día a día, se consideran \textbf{inembargables}.

La Corte desarrolla nuevamente la tesis de la prevalencia de la igualdad religiosa, del tratamiento igualitario y de la no preferencia. El sostenimiento del culto del que habla el artículo 2 debe interpretarse con un criterio restrictivo, limitado al aspecto económico institucional. Esta es la jurisprudencia actual, que profundiza lo que se venía diciendo desde 1994.

\subsection{Exenciones impositivas y ejecución de bienes}

La Iglesia tiene exenciones a la tributación en la mayoría de los casos, sobre todo en los impuestos de alcance nacional; cada provincia regula sus propios impuestos. Se trata de una desgravación impositiva que forma parte del sostenimiento del culto, y que se vincula, más allá del sostenimiento religioso, con la idea de que no hay ánimo de lucro en las actividades que despliega la Iglesia.

La Iglesia es dueña de importantes propiedades dentro del país. En la medida en que no estén afectadas al culto católico diario, esas propiedades son susceptibles de ejecución judicial si hubiera una condena a pagar contra la persona jurídica. Puede ser demandada, por ejemplo, la arquidiócesis, que es una unidad administrativa y jurisdiccional independiente dentro de la persona jurídica Iglesia Católica.

\begin{important}
La Iglesia Católica es persona jurídica pública, con concordato de rango de tratado, bienes afectados al culto inembargables y exenciones impositivas nacionales. Los bienes no afectados al culto diario pueden ser ejecutados, y la arquidiócesis puede ser demandada.
\end{important}

\section{Cuadro Cornell del capítulo}

\begin{longtable}{
    >{\raggedright\arraybackslash}p{0.27\textwidth}
    !{\vrule width 0.4pt}
    >{\raggedright\arraybackslash}p{\dimexpr0.73\textwidth-6\tabcolsep-0.4pt\relax}
}
\toprule
\textbf{Preguntas / palabras clave} & \textbf{Notas / respuestas} \\
\midrule
\endfirsthead
\toprule
\textbf{Preguntas / palabras clave} & \textbf{Notas / respuestas} \\
\midrule
\endhead
\bottomrule
\endfoot
\bottomrule
\endlastfoot

\textbf{¿Por qué se eliminó en 1994 el requisito de pertenecer al culto católico?} &
Era uno de los puntos principales de la temática religiosa y se suprimió directamente en 1994. Afianzó la idea de que el culto no debe influir en la gestión política e institucional del Estado. Dejó de ser requisito habilitante, y la Convención Constituyente reafirmó la igualdad religiosa y el trato igualitario de los cultos oficialmente reconocidos (registro de cultos). \\[0.6em]

\textbf{¿Qué cuestionó la demanda contra Salta: la norma o su aplicación?} &
Una asociación demandó a la provincia por la educación religiosa en los colegios. La norma provincial garantizaba la enseñanza de la religión que cada alumno profese; se atacó su aplicación práctica y no el texto. La Corte falló casi por unanimidad, con disidencia parcial del juez Rosatti. \\[0.6em]

\textbf{¿Qué alcance tienen el artículo 2 y el artículo 14, y qué es el principio de neutralidad?} &
El artículo 2 (``sostiene el culto católico'') se limita a un sostén institucional y económico, con interpretación restrictiva. El artículo 14 asegura la libertad de culto, y el Estado debe practicar igualdad, no discriminación y no preferencia. La neutralidad parte de una sentencia del tribunal alemán y es común en Europa. \\[0.6em]

\textbf{¿Puede ser inconstitucional la aplicación de una norma de texto válido?} &
Sí. Se probó que la religión se daba en todos los cursos, que solo se enseñaba catecismo católico, que había un formulario para declarar el culto y que los no católicos perdían horas sin alternativa. La Corte declaró inválida la práctica por violar la neutralidad en la educación; la provincia conserva su competencia sobre la educación básica, siempre con trato igualitario. \\[0.6em]

\textbf{¿Qué estatus jurídico tiene la Iglesia Católica?} &
Es persona jurídica pública (Código Civil y Comercial) y tiene un concordato con rango de tratado internacional. Sus bienes afectados al culto son inembargables y tiene exenciones impositivas nacionales. Los bienes no afectados al culto diario son ejecutables, y la arquidiócesis puede ser demandada. \\

\midrule
\multicolumn{2}{p{\dimexpr\textwidth-4\tabcolsep\relax}}{%
\textbf{Resumen:} la supresión del requisito religioso en 1994 reafirmó la igualdad de los cultos. La Corte interpreta el artículo 2 como un sostén solo institucional y económico, y el artículo 14 exige neutralidad. En el caso Salta, una norma de texto válido fue declarada inválida en su aplicación práctica porque, en los hechos, solo se enseñaba la religión católica, se obligaba a exteriorizar el culto y los alumnos no católicos perdían horas sin alternativa. La Iglesia Católica es persona jurídica pública, con bienes inembargables afectados al culto y exenciones impositivas nacionales.
} \\
\end{longtable}

% ==================================================================
% CAPÍTULO 2
% ==================================================================
\chapter[Igualdad y categorías sospechosas]{Igualdad, categorías sospechosas y carga de la prueba}

\section{Categorías sospechosas e inversión de la carga de la prueba}

\subsection{La presunción de legitimidad de las normas}

Cuando existen regulaciones que, en los hechos o en la letra de la norma, implican un trato preferente por motivos religiosos, entran en el cono de sospecha de las \textbf{categorías sospechosas} de las que habla la Corte.

\begin{definition}{Categorías sospechosas}
Normas que hacen un distingo fundado en temas objetivos como la religión, el género, la edad o la preferencia política. En ellas es el órgano estatal el que tiene que demostrar que la norma es válida: se invierte la presunción de legitimidad que alcanza a cualquier norma regularmente sancionada.
\end{definition}

Toda norma regularmente sancionada goza de una presunción de legitimidad de tipo \textit{iuris tantum}. Por ejemplo, quien plantea una medida cautelar suspendiendo una norma contra el Estado nacional debe rebatir esa presunción. Una de las formas en que esa presunción se modifica es que la norma tenga un trato diferenciado por cuestiones objetivas de composición religiosa, de género, etaria o de preferencia política. Ello pone a la norma el cartel de ``esta norma puede ser inconstitucional'', y quien la sancionó es quien debe demostrar su legalidad: se invierte la carga de la prueba. La Corte lo ha desarrollado en ese sentido.

\begin{table}[H]
\centering
\caption{Norma común frente a norma en categoría sospechosa}
\small
\begin{tabularx}{\textwidth}{
    >{\raggedright\arraybackslash}p{0.22\textwidth}
    >{\raggedright\arraybackslash}X
    >{\raggedright\arraybackslash}X
}
\toprule
\textbf{Criterio} & \textbf{Norma común} & \textbf{Categoría sospechosa} \\
\midrule
\textbf{Presunción} & Se presume legítima si fue sancionada por los mecanismos legales. & Se sospecha de su invalidez hasta que se demuestre lo contrario. \\
\textbf{Quién carga con la prueba} & Quien intenta invalidarla debe probar la invalidez. & El órgano estatal debe demostrar que la norma es válida. \\
\textbf{Cuándo se activa} & Siempre, como regla general. & Cuando distingue por religión, género, edad o preferencia política. \\
\textbf{Efecto práctico} & Se aplica mientras no se declare inválida. & No deja de aplicarse, pero el Estado debe explicar su razonabilidad. \\
\bottomrule
\end{tabularx}
\end{table}

\begin{important}
Las normas gozan de presunción de legitimidad, pero cuando distinguen por religión, género o edad se invierte la carga de la prueba: el Estado debe demostrar que la norma es razonable y válida.
\end{important}

\section{Ejemplo cotidiano: el límite de edad docente}

Cuando una norma hace una diferenciación que genera un gravamen o perjuicio para una determinada característica de las personas humanas, entra en el cono de sospecha y se invierte la carga de la prueba. La norma se sospecha inválida, a priori, hasta que se demuestre lo contrario en el proceso judicial. No es que la norma no se aplique: quien invoca la invalidez puede invocar este precepto, y el Estado debe explicar por qué la norma es razonable.

\begin{example}{El límite de 70 años para dar clases}
No se pueden impartir clases en la Universidad de Buenos Aires cumplidos los 70 años, salvo que exista un estatus como profesor emérito u otro régimen que permita seguir dictando clases. La norma hace perder un derecho por el paso del tiempo, por una cuestión de edad. Por ello podría reputarse discriminatoria y plantearse su ilegitimidad, de modo que entra en categoría sospechosa.

Si un profesor quiere cuestionarla, debe plantearlo, y el Estado, al contestar la demanda, debe explicar las razones y la motivación de la norma. Debe justificar que es razonable, más allá de las aptitudes y el estado físico de cada profesor en particular, por:
\begin{itemize}
    \item la necesidad de generar una renovación del cuerpo docente;
    \item la necesidad de evitar perpetuarse en los cargos;
    \item una posible declinación en la intensidad de atención al impartir clases.
\end{itemize}
Como la regla se aplica de forma objetiva, siempre hay una válvula de escape para casos concretos: figuras como la del profesor emérito, que funcionan como reconocimiento y convalidación de que la persona puede seguir dando clases.
\end{example}

\begin{important}
Ante una demanda contra una norma por edad, el Estado debe explicar su razonabilidad: renovación del cuerpo docente, evitar la perpetuación en los cargos, posible declinación en la atención y válvulas de escape como la figura del profesor emérito.
\end{important}

\section{Nacionalidad y acceso a cargos públicos}

\subsection{El caso de los secretarios de juzgados nacionales}

Se mencionó una demanda por la invalidez de una disposición de la justicia nacional que determinaba que los secretarios de los juzgados nacionales debían ser de nacionalidad argentina. Fue impugnada por una persona de apellido alemán. El caso se llama \textbf{Gottschau}. La Corte le dio la razón, entendiendo que la diferencia fundada en la nacionalidad para un cargo como el de secretario de un juzgado de las cámaras nacionales no ameritaba esa restricción, y que debía declararse inconstitucional en aras de evitar conductas discriminatorias.

% TODO: contenido incompleto o ambiguo en las notas originales (el nombre «Gottschau» se escucha deformado en la grabación y el docente añade un comentario sobre su familia que no se comprende)

\subsection{Cargos distintos, análisis distinto}

\begin{example}{El jefe del Estado Mayor del Ejército}
Una cosa es el cargo de secretario de un juzgado civil, laboral o de cualquiera de las cámaras nacionales, y otra el de jefe del Estado Mayor del Ejército. La obligatoriedad de la nacionalidad argentina de nacimiento, o de ser hijo de naturales del país, podría reputarse inconstitucional también en este caso. Hay que apreciar la naturaleza del cargo y sus condiciones, y pensar si el cargo de jefe del ejército de un país puede restringirse razonablemente a nacionales. El docente aventura, como opinión propia, que probablemente la Corte tendría una posición diferente.
\end{example}

Algo similar ocurre con la figura de presidente o presidenta de la Nación, que debe ser ciudadano nativo o hijo de ciudadanos argentinos.

\begin{quote}
\textbf{Art. 89, Constitución Nacional (norma aludida; sin cita expresa del número).} Para ser elegido presidente o vicepresidente de la Nación, se requiere haber nacido en el territorio argentino, o ser hijo de ciudadano nativo, habiendo nacido en país extranjero; y las demás calidades exigidas para ser elegido senador.
\end{quote}

\subsection{Comparación con Estados Unidos y un caso local}

\begin{example}{La impugnación contra Barack Obama}
En la Constitución de los Estados Unidos ni siquiera está la posibilidad de que el presidente sea hijo de ciudadanos estadounidenses: debe haber nacido en territorio norteamericano. Se recordó una impugnación, calificada de delirante y sin fundamento, que sostenía que Obama había nacido en Hawái, en los años 60, y que Hawái no formaba parte de los Estados Unidos en ese momento, por lo que no cumpliría el requisito.
\end{example}

% TODO: contenido incompleto o ambiguo en las notas originales (el razonamiento del docente sobre la fecha de incorporación de Hawái es confuso)

\begin{example}{El origen de De Narváez}
De Narváez había nacido en Colombia y no era de padres argentinos. Cuando era una figura política en su momento de mayor protagonismo, estaba en discusión si podía ser candidato o no, y era un obstáculo insalvable para él.
\end{example}

Las cláusulas que puedan reputarse discriminatorias o sospechosas deben analizarse muy en detalle, según el ámbito en el que se define la característica que deja atrás a unos y favorece a otros. Por ejemplo, en los requisitos para un trabajo importa mucho el grado de responsabilidad política institucional de ese trabajo.

\begin{table}[H]
\centering
\caption{La exigencia de nacionalidad según el cargo (ejemplos del docente)}
\small
\begin{tabularx}{\textwidth}{
    >{\raggedright\arraybackslash}p{0.22\textwidth}
    >{\raggedright\arraybackslash}p{0.22\textwidth}
    >{\raggedright\arraybackslash}X
}
\toprule
\textbf{Cargo} & \textbf{Requisito de nacionalidad} & \textbf{Análisis planteado} \\
\midrule
\textbf{Secretario de juzgado nacional} & Nacionalidad argentina. & La Corte lo declaró inconstitucional (caso Gottschau): la restricción no se justificaba por la jerarquía del cargo. \\
\textbf{Jefe del Estado Mayor del Ejército} & Nativo o hijo de nativo. & Hay que apreciar la naturaleza del cargo; probablemente la Corte tendría una posición diferente. \\
\textbf{Presidente de la Nación (Argentina)} & Nativo o hijo de ciudadano nativo. & Requisito constitucional expreso, por el grado de responsabilidad política institucional. \\
\textbf{Presidente (Estados Unidos)} & Nacido en territorio norteamericano. & Discusión sobre el origen de Obama (sin fundamento). \\
\bottomrule
\end{tabularx}
\end{table}

\begin{important}
La exigencia de nacionalidad se analiza según la naturaleza y la responsabilidad política institucional del cargo: lo que no se justifica para un secretario de juzgado puede ser razonable para otros cargos.
\end{important}

\section{La edad de los jueces: doctrinas Fayt y Schiffrin}

\subsection{La edad como criterio: de los trabajos a la magistratura}

Cuando el criterio es la edad también deben tenerse en cuenta la exigencia de determinadas labores y el aspecto, que no es menor, de la renovación académica. Un ejemplo es el de los 75 años, en que los jueces cesan sus funciones. Es la única causal que el régimen federal establece para el cese de un juez que, mientras dure su buena conducta, debe ser independiente, y una condición de ello es su inamovilidad en el cargo.

\begin{quote}
\textbf{Art. 110, Constitución Nacional (norma aludida: ``mientras dure su buena conducta'').} Los jueces de la Corte Suprema y de los tribunales inferiores de la Nación conservarán sus empleos mientras dure su buena conducta, y recibirán por sus servicios una compensación que determinará la ley, y que no podrá ser disminuida en manera alguna, mientras permaneciesen en sus funciones.
\end{quote}

\begin{quote}
\textbf{Art. 99, inc. 4 (párrafos finales), Constitución Nacional (norma aludida: límite de 75 años).} [\ldots] Un nuevo nombramiento, precedido de igual acuerdo, será necesario para mantener en el cargo a cualquiera de esos magistrados, una vez que cumplan la edad de setenta y cinco años. Todos los nombramientos de magistrados cuya edad sea la indicada o mayor se harán por cinco años, y podrán ser repetidos indefinidamente, por el mismo trámite.
\end{quote}

\subsection{La doctrina Fayt y la doctrina Schiffrin}

La doctrina Fayt rigió durante mucho tiempo; luego vino la doctrina Schiffrin, un fallo posterior. En ella la Corte dijo que la restricción de los 75 años para desempeñar el cargo de juez era \textbf{razonable}. Además, estaba incluida dentro de las posibilidades que tenían los constituyentes para regular, porque el punto de la designación de los jueces había sido considerado un punto susceptible de reforma.

% TODO: contenido incompleto o ambiguo en las notas originales (no se detalla el contenido de la doctrina Fayt)

\begin{important}
La Corte (doctrina Schiffrin, posterior a la doctrina Fayt) consideró razonable el límite de 75 años para los jueces y comprendido dentro de lo que los constituyentes podían regular.
\end{important}

\section{Género, igualdad y proporcionalidad}

\subsection{Casos del ámbito privado: heladerías y colectivos}

Se mencionaron dos precedentes del ámbito privado:
\begin{itemize}
    \item el fallo \textbf{Freddo}, sobre la contratación de mujeres en las heladerías;
    \item un fallo sobre conductoras de colectivos de media y corta distancia, que corresponde al caso \textbf{Sisnero}.
\end{itemize}

% TODO: contenido incompleto o ambiguo en las notas originales (la alumna ubica el fallo sobre conductores en Jujuy; la nota editorial indica que tramitó en la provincia de Salta)

\subsection{El género como punto central: la Convención contra la discriminación de la mujer}

Uno de los puntos centrales se da con el tema de género, donde históricamente ha habido un trato discriminatorio e inequitativo. Esto se afianza con un tratado suscripto al efecto: la \textbf{Convención sobre la Eliminación de Todas las Formas de Discriminación contra la Mujer (CEDAW)}. Las notas indican que tiene jerarquía constitucional conforme al art. 75, inc. 22, de la Constitución Nacional.

\subsection{Legalidad y proporcionalidad de las diferencias}

Para las diferencias fundadas en género se señalan dos exigencias:
\begin{itemize}
    \item \textbf{Legalidad:} las diferencias deben estar taxativamente previstas de antemano y bien detalladas, para saber quiénes pueden aplicar a un determinado estatus o no.
    \item \textbf{Proporcionalidad:} debe haber proporcionalidad entre las exigencias de las características que se invoquen y el desempeño del cargo.
\end{itemize}

En el ejemplo de los colectivos, la Corte sostuvo que las habilidades y capacidades no se miden por quién maneja mejor, sino a veces por cuestiones como la resistencia física. Pero afirmó que la actividad de conductor, más allá de demandar cierto esfuerzo, no ponía en un pie de desigualdad a hombres y mujeres respecto de sus capacidades y aptitudes físicas.

\begin{important}
Las diferencias por género deben estar previstas de antemano, detalladas y ser proporcionadas a las exigencias reales del cargo; en el caso de los conductores, la Corte descartó una desigualdad de aptitudes físicas entre hombres y mujeres.
\end{important}

\section{Cuadro Cornell del capítulo}

\begin{longtable}{
    >{\raggedright\arraybackslash}p{0.27\textwidth}
    !{\vrule width 0.4pt}
    >{\raggedright\arraybackslash}p{\dimexpr0.73\textwidth-6\tabcolsep-0.4pt\relax}
}
\toprule
\textbf{Preguntas / palabras clave} & \textbf{Notas / respuestas} \\
\midrule
\endfirsthead
\toprule
\textbf{Preguntas / palabras clave} & \textbf{Notas / respuestas} \\
\midrule
\endhead
\bottomrule
\endfoot
\bottomrule
\endlastfoot

\textbf{¿Quién debe probar la validez de una norma que distingue por religión, género o edad?} &
Las normas comunes tienen presunción de legitimidad (\textit{iuris tantum}), y quien las impugna debe probar su invalidez. Las normas que distinguen por religión, género, edad o preferencia política son categorías sospechosas: se invierte la carga y el Estado debe demostrar su validez. La norma no deja de aplicarse. \\[0.6em]

\textbf{¿Es discriminatorio retirar a un docente por su edad y cómo defiende el Estado esa norma?} &
El límite de 70 años en la UBA se sospecha discriminatorio y, al contestar la demanda, el Estado debe justificarlo: renovación del cuerpo docente, evitar la perpetuación en los cargos y posible declinación en la atención. La válvula de escape es la figura del profesor emérito. \\[0.6em]

\textbf{¿Cuándo es razonable exigir la nacionalidad argentina para un cargo?} &
Se analiza el grado de responsabilidad política institucional del cargo. En Gottschau, la exigencia para secretarios de juzgado fue declarada inconstitucional. Para jefe del Estado Mayor, el docente supone que la Corte tendría una posición distinta; para presidente rige el art. 89 (nativo o hijo de nativo). \\[0.6em]

\textbf{¿Es razonable que los jueces cesen a los 75 años?} &
Los 75 años son la única causal de cese del juez, que conserva su cargo mientras dure su buena conducta. En la doctrina Schiffrin, posterior a la doctrina Fayt, la Corte consideró razonable el límite y comprendido en lo que los constituyentes podían regular. \\[0.6em]

\textbf{¿Qué exigencias rigen para las diferencias por género?} &
Las diferencias deben estar previstas de antemano y detalladas (legalidad) y las exigencias deben ser proporcionales al desempeño del cargo. La CEDAW afianza la necesidad de evitar el trato discriminatorio. La Corte, en el caso de conductores, sostuvo que conducir no genera desigualdad de aptitudes físicas entre hombres y mujeres. \\

\midrule
\multicolumn{2}{p{\dimexpr\textwidth-4\tabcolsep\relax}}{%
\textbf{Resumen:} las normas gozan de presunción de legitimidad, pero cuando distinguen por religión, género, edad o preferencia política entran en categoría sospechosa y el Estado debe demostrar su razonabilidad. Lo que se justifica para un cargo puede no justificarse para otro: la nacionalidad y la edad se analizan según la naturaleza y la responsabilidad del cargo. En materia de género, las diferencias deben estar previstas de antemano y ser proporcionadas.
} \\
\end{longtable}

% ==================================================================
% CAPÍTULO 3
% ==================================================================
\chapter[Defensa del consumidor]{Defensa del consumidor y poder sancionatorio}

\section{Poder sancionatorio en Defensa del Consumidor}

\subsection{Quién sanciona y con qué sanciones}

La Secretaría de Comercio, por intermedio de la Subsecretaría de Defensa del Consumidor, tiene la posibilidad de sancionar a empresas, al igual que los organismos provinciales. Puede hacerlo cuando, en investigaciones de oficio o a partir de denuncias, se corroboran conductas abusivas o perjudiciales para el consumidor.

\begin{formula}{Sanciones de la Ley 24.240 de Defensa del Consumidor (según lo expuesto)}
\begin{itemize}
    \item \textbf{Multa:} es la sanción más común y de contenido pecuniario. Puede llegar hasta \textbf{2100 canastas básicas}, de las que publica el INDEC.
    \item \textbf{Inhabilitación} de los locales comerciales.
    \item \textbf{Exclusión} de los registros de proveedores del Estado.
\end{itemize}
\end{formula}

% TODO: contenido incompleto o ambiguo en las notas originales (el docente describe las sanciones sin citar el número de artículo)

\subsection{No se exige acreditar un perjuicio concreto}

En estas investigaciones muchas veces no se necesita acreditar de forma fehaciente un perjuicio concreto y verificable. La propia concreción de ciertas conductas puede dar lugar a la sanción, por ejemplo:
\begin{itemize}
    \item ofertas publicitarias inexactas;
    \item información que no es veraz;
    \item elementos que puedan llevar a confusión al consumidor;
    \item la introducción masiva, en los contratos de empresas que se relacionan con consumidores, de cláusulas abusivas, como las que invierten la carga de la prueba en materia de responsabilidad.
\end{itemize}

Esta sanción debe llevarse a cabo en el marco de un procedimiento sancionatorio regulado en la Ley 24.240, ante la autoridad administrativa, donde debe asegurarse el \textbf{debido proceso}. Es un aspecto clave para ejercer legítimamente el poder punitivo. Con este poder se busca desalentar conductas que tienen cierto grado de masividad, es decir, conductas desplegadas en la actuación general de las empresas, donde puede haber una cantidad indeterminada de consumidores perjudicados.

\begin{important}
El poder punitivo busca desalentar conductas masivas de las empresas; por eso muchas infracciones se configuran por el solo hecho de publicitar información inexacta o de incluir cláusulas abusivas, sin probar un daño concreto.
\end{important}

\section{Ofertas publicitarias e incumplimiento}

Entre los tipos sancionatorios figuran la información inexacta y el incumplimiento de oferta. Una oferta comercial indeterminada, dirigida en medios masivos a una pluralidad de consumidores, debe tener especificadas y debidamente comunicadas, en la misma publicidad, sus condiciones. Por ejemplo, debe informar cuántas son las unidades en venta, para no frustrar las expectativas de los consumidores.

\begin{example}{La oferta de carne y de televisores}
Es típico que supermercados o comercios masivos anuncien, en todos los locales del radio de la Capital Federal, ``venta de carne este fin de semana a 10.000 pesos el kilo'' o televisores a determinado precio. Una oferta así genera una expectativa muy grande, y ese fin de semana se presentarán miles de consumidores. Si los consumidores se presentan sin aviso previo y no obtienen el producto, hay un incumplimiento de la oferta.

Es distinto si la publicidad dice ``hasta un máximo de 3.000 kg de carne'' o ``300 televisores por sucursal''. Esa información da una idea de la posible dimensión de la oferta y hace saber a los consumidores que existe la posibilidad de que uno vaya y no tenga, o vea frustrada su expectativa. No es lo mismo que si no hay ningún aviso al respecto.
\end{example}

Estos incumplimientos son, además, ardides publicitarios para apalancar ventas masivas y mover la atención del consumidor. La infracción no necesita acreditar fehacientemente un perjuicio sobre un número concreto de consumidores: en muchos casos, el propio incumplimiento exterioriza una eventual violación a la ley.

\begin{table}[H]
\centering
\caption{Oferta con y sin información sobre el alcance}
\small
\begin{tabularx}{\textwidth}{
    >{\raggedright\arraybackslash}p{0.23\textwidth}
    >{\raggedright\arraybackslash}X
    >{\raggedright\arraybackslash}X
}
\toprule
\textbf{Situación} & \textbf{Información en la publicidad} & \textbf{Efecto sobre el consumidor} \\
\midrule
\textbf{Oferta sin límites informados} & ``Carne a 10.000 pesos el kilo'', sin indicar cantidades. & Expectativa masiva; si no hay producto, hay incumplimiento de la oferta. \\
\textbf{Oferta con alcance informado} & ``Hasta un máximo de 3.000 kg'' o ``300 televisores por sucursal''. & El consumidor sabe que puede no alcanzar; no se frustra su expectativa sin aviso. \\
\bottomrule
\end{tabularx}
\end{table}

\begin{important}
Una oferta masiva debe comunicar sus condiciones y alcance en la misma publicidad; el incumplimiento de la oferta exterioriza la violación a la ley sin necesidad de probar un perjuicio concreto.
\end{important}

\section{Derecho administrativo sancionatorio y principios penales}

\subsection{Garantías del procedimiento sancionatorio}

El procedimiento sancionatorio debe asegurar el resguardo de las garantías del debido proceso a la empresa investigada:
\begin{itemize}
    \item una oportunidad de presentar su descargo;
    \item una oportunidad de producir prueba, y que esa prueba sea considerada y analizada al decidir sobre la sanción;
    \item una imputación concreta al momento de ser citada al procedimiento, que permita un ejercicio efectivo del derecho de defensa: debe decirse qué tipo de conducta se le imputa.
\end{itemize}

\begin{quote}
\textbf{Art. 18, Constitución Nacional (norma aludida: principios penales, inocencia y defensa).} Ningún habitante de la Nación puede ser penado sin juicio previo fundado en ley anterior al hecho del proceso, ni juzgado por comisiones especiales, o sacado de los jueces designados por la ley antes del hecho de la causa. [\ldots] Es inviolable la defensa en juicio de la persona y de los derechos.
\end{quote}

\subsection{Un derecho administrativo sancionatorio con principios del derecho penal}

Estamos en el ámbito del \textbf{derecho administrativo sancionatorio}. Este derecho no es un subconjunto o una ``clase B'' del derecho penal, pues tiene sus propias características. Pero como el objetivo final es el poder punitivo, que además debe ser proporcionado, legítimo y concordante con la finalidad de la norma, el Estado debe asegurar que se cubran los principios del derecho penal material o sustantivo.

\subsection{Tipicidad}

\begin{definition}{Tipicidad}
El tipo sancionatorio debe estar especificado en la ley y debe demostrarse que la conducta está en el presupuesto fijado por ella. Al igual que en el derecho penal, en el derecho sancionatorio debe ser clara y evidente.
\end{definition}

Por ejemplo, para sancionar a quien incumple un régimen promocional difundido en medios masivos a una pluralidad de consumidores, debe demostrarse que la oferta no fue cumplida y que no hubo debida información al destinar la promoción a los consumidores.

\subsection{Ley penal más benigna e irretroactividad}

Si la ley es dejada sin efecto posteriormente, en principio el órgano sancionatorio no puede mantener la sanción, porque la ley quedó derogada. Si un tipo de infracción queda sin efecto, el Estado no puede mantener el ejercicio del monopolio del poder punitivo, porque el legislador ya consideró que eso no es un incumplimiento de la ley. Rigen así la tipicidad, la aplicación de la ley penal más benigna y la irretroactividad de la norma: la norma debe estar plenamente vigente, en sus elementos sustanciales, al momento de la comisión del hecho que se endilga.

\subsection{\textit{Non bis in idem}}

\begin{definition}{\textit{Non bis in idem}}
No puede haber dos sanciones por un mismo hecho. Pueden coexistir una infracción en el derecho penal sustantivo y una violación de la ley de defensa del consumidor, pero no dos sanciones dentro del mismo órgano o del mismo ámbito punitivo.
\end{definition}

\subsection{Debido proceso y resolución fundada}

El procedimiento sancionatorio debe cumplir todas las garantías del debido proceso: debe imputarse la infracción que se sospecha, la empresa debe poder defenderse y producir prueba, y la resolución sancionatoria debe estar debidamente fundamentada en normas vigentes y aplicables a esa persona jurídica. Aun cuando no sea una materia de derecho penal, el ámbito sancionatorio del Estado debe cumplir, cuanto menos, con los principios generales del derecho penal. Según la jurisprudencia aplicable, el derecho administrativo sancionatorio debe cumplir los principios esenciales del derecho penal, aunque no haya una equivalencia con el proceso penal.

\begin{table}[H]
\centering
\caption{Principios del derecho penal exigibles al derecho administrativo sancionatorio}
\small
\begin{tabularx}{\textwidth}{
    >{\raggedright\arraybackslash}p{0.25\textwidth}
    >{\raggedright\arraybackslash}X
}
\toprule
\textbf{Principio} & \textbf{Contenido según el docente} \\
\midrule
\textbf{Tipicidad} & La infracción debe estar especificada en la ley y demostrarse que la conducta está en el presupuesto fijado por ella; debe ser clara y evidente. \\
\textbf{Ley penal más benigna} & Si la ley es dejada sin efecto, el órgano no puede mantener la sanción, porque el legislador ya consideró que no es un incumplimiento. \\
\textbf{Irretroactividad} & La norma debe estar vigente en sus elementos sustanciales al momento del hecho que se endilga. \\
\textbf{\textit{Non bis in idem}} & No puede haber dos sanciones por el mismo hecho dentro del mismo ámbito punitivo (sí pueden coexistir en campos jurídicos distintos). \\
\textbf{Debido proceso} & Imputación concreta, descargo, producción y consideración de la prueba, y resolución fundada en normas vigentes y aplicables. \\
\bottomrule
\end{tabularx}
\end{table}

\begin{important}
Aunque no sea derecho penal, cuando el Estado despliega poder punitivo debe respetar los principios esenciales del derecho penal sustantivo y las garantías del debido proceso.
\end{important}

\section{Arbitrajes, garantías mínimas y servicios públicos}

\subsection{Arbitrajes en las relaciones de consumo}

Para organizar el mercado y proteger al consumidor, la Secretaría regula los \textbf{arbitrajes en relaciones de consumo}. Prevé ámbitos en los que capacita árbitros que median entre dos partes en conflicto. Ya no se trata de la Administración que persigue a un incumplidor de la Ley 24.240: hay dos partes en conflicto y el árbitro es una instancia de resolución de conflictos. Lo que resuelva el árbitro se puede impugnar luego ante la justicia. En cuanto a la posibilidad de multa, llega hasta 2100 canastas básicas, informadas por el INDEC.

\subsection{Garantías: un piso de orden público}

La ley establece un piso de garantía para los bienes que se comercializan. Se trata de una norma de orden público, que predomina incluso sobre el acuerdo de partes.

\subsection{Servicios públicos domiciliarios}

Algunas cuestiones referidas al derecho del consumidor de servicios públicos domiciliarios se complementan con lo que prevé cada marco regulatorio, por ejemplo el del transporte y distribución del gas o de la energía. Algunos pocos conceptos del servicio público domiciliario también están en la Ley 24.240.

\begin{formula}{Reglas mencionadas de la Ley 24.240 (según lo expuesto)}
\begin{itemize}
    \item \textbf{Garantía legal mínima:} 6 meses para bienes nuevos y 3 meses para bienes usados. Es un piso de orden público que prevalece sobre el acuerdo de las partes.
    \item \textbf{Servicios públicos domiciliarios:} si la facturación supera en más de un 75\% el consumo promedio para esa época del año, el usuario puede no pagar y pedir la revisión.
\end{itemize}
\end{formula}

% TODO: contenido incompleto o ambiguo en las notas originales (el docente no cita el número de artículo de estas reglas)

\begin{example}{La tarifa domiciliaria que aumenta más de un 75\%}
Cuando la tarifa domiciliaria de gas, electricidad o agua aumenta más de un 75\% respecto del consumo promedio para esa época del año (no del precio, que aumenta todo el tiempo), el usuario tiene derecho a no pagarla y pedir su revisión, como garantía para evitar cualquier lesión, porque se presume que puede ser una tarifa mal calculada. La boleta no se paga y se puede discutir; si no se demuestra que estaba mal calculada, el servicio deberá pagarse.
\end{example}

\begin{important}
El arbitraje de consumo resuelve conflictos entre dos partes y puede impugnarse ante la justicia; las garantías legales (6 y 3 meses) son un piso de orden público; y el aumento de tarifa superior al 75\% permite no pagar y pedir revisión.
\end{important}

\section{Cuadro Cornell del capítulo}

\begin{longtable}{
    >{\raggedright\arraybackslash}p{0.27\textwidth}
    !{\vrule width 0.4pt}
    >{\raggedright\arraybackslash}p{\dimexpr0.73\textwidth-6\tabcolsep-0.4pt\relax}
}
\toprule
\textbf{Preguntas / palabras clave} & \textbf{Notas / respuestas} \\
\midrule
\endfirsthead
\toprule
\textbf{Preguntas / palabras clave} & \textbf{Notas / respuestas} \\
\midrule
\endhead
\bottomrule
\endfoot
\bottomrule
\endlastfoot

\textbf{¿Cómo y por qué sanciona el Estado a las empresas que perjudican al consumidor?} &
La Secretaría de Comercio (Subsecretaría de Defensa del Consumidor) y los organismos provinciales investigan de oficio o por denuncia. Las sanciones son: multa de hasta 2100 canastas básicas (INDEC), inhabilitación de locales y exclusión de registros de proveedores del Estado. No siempre hace falta acreditar un perjuicio concreto: basta, por ejemplo, la publicidad inexacta o las cláusulas abusivas. El poder punitivo busca desalentar conductas masivas. \\[0.6em]

\textbf{¿Qué debe informar una oferta masiva para no incumplir?} &
Sus condiciones y alcance en la misma publicidad (por ejemplo, ``hasta 3.000 kg'' o ``300 televisores por sucursal''). Sin esa información, que los consumidores se presenten y no obtengan el producto es un incumplimiento de la oferta, que exterioriza la violación a la ley. \\[0.6em]

\textbf{¿Qué garantías del derecho penal rigen cuando sanciona la Administración?} &
Tipicidad, ley penal más benigna, irretroactividad, \textit{non bis in idem} y debido proceso (imputación concreta, descargo, prueba y resolución fundada). El art. 18 CN dice que nadie puede ser penado sin juicio previo fundado en ley anterior al hecho y que es inviolable la defensa en juicio. \\[0.6em]

\textbf{¿Qué es el arbitraje de consumo y cuál es su relación con la justicia?} &
Es una instancia de resolución de conflictos entre dos partes, con árbitros capacitados por la Secretaría. Lo que resuelva el árbitro se puede impugnar ante la justicia. \\[0.6em]

\textbf{¿Qué derechos mínimos del consumidor no puede quitar ningún contrato?} &
La garantía mínima de 6 meses (bienes nuevos) y 3 meses (usados) es de orden público y predomina sobre el acuerdo de partes. Si la tarifa domiciliaria supera en más de un 75\% el consumo promedio de esa época del año, se puede no pagar y pedir revisión. \\

\midrule
\multicolumn{2}{p{\dimexpr\textwidth-4\tabcolsep\relax}}{%
\textbf{Resumen:} el Estado puede sancionar a las empresas con multas de hasta 2100 canastas básicas, inhabilitación o exclusión de registros, muchas veces sin probar un perjuicio concreto, pero debe respetar el debido proceso y los principios del derecho penal (tipicidad, ley más benigna, irretroactividad, \textit{non bis in idem}). Las ofertas masivas deben informar su alcance. Quedan reglas mínimas de orden público: garantía de 6 y 3 meses y posibilidad de no pagar y pedir revisión ante aumentos superiores al 75\%.
} \\
\end{longtable}

% ==================================================================
% CAPÍTULO 4
% ==================================================================
\chapter[Solve et repete y lealtad comercial]{\textit{Solve et repete} y lealtad comercial}

\section{El \textit{solve et repete} y el principio de inocencia}

\subsection{La reforma de 2014 y el \textit{solve et repete}}

En su última reforma, del año 2014, la ley establecía que las sanciones impuestas por la Secretaría de Comercio, en el área de Defensa del Consumidor, debían pagarse como forma de poder acceder a la revisión judicial.

\begin{definition}{\textit{Solve et repete}}
Régimen por el cual, para cuestionar una multa, primero hay que pagarla como presupuesto de acceso a la sede judicial.
\end{definition}

Era un precepto bastante general en alguna época. Muchas normas establecían que, aunque la sanción fuera impuesta por un órgano administrativo y todo sujeto tuviera derecho a cuestionarla en sede judicial, se iba a la justicia pagando la multa como presupuesto de acceso.

\subsection{Problemas constitucionales}

En clase se identificaron dos aspectos afectados:
\begin{itemize}
    \item \textbf{Propiedad:} se agrede el patrimonio, pues hay que destinar parte de él a pagar una multa para recién luego cuestionarla.
    \item \textbf{Principio de inocencia:} si hay que pagar al principio, la pena ya está aplicada. Aunque el dinero sea fungible y pueda retornarse después, hay una presunción de culpabilidad y una violación del principio de inocencia. Por más que se le pongan los rótulos que se quiera, al final del día es un órgano administrativo el que agrede el patrimonio.
\end{itemize}

\subsection{La solución de la Ley de Procedimientos Administrativos}

Este punto está zanjado con las modificaciones a la Ley de Procedimientos Administrativos, que establece, de forma genérica, la prohibición de dispositivos que prevean el \textit{solve et repete}, es decir, la obligación de pagar las multas que imponga la Administración. Es una autorrestricción del propio Estado, que el docente considera correcta: a partir de ahora no puede cobrar la multa. Esto vale para la Ley de Defensa del Consumidor y también, por ejemplo, para las multas millonarias en defensa de la competencia o en el decreto-ley de lealtad comercial.

% TODO: contenido incompleto o ambiguo en las notas originales (las fechas y alcances de la reforma y de la prohibición no se precisan; el docente menciona que la prohibición rige ``desde hace 2 años'')

\begin{table}[H]
\centering
\caption{El régimen del \textit{solve et repete}}
\small
\begin{tabularx}{\textwidth}{
    >{\raggedright\arraybackslash}p{0.22\textwidth}
    >{\raggedright\arraybackslash}X
    >{\raggedright\arraybackslash}X
}
\toprule
\textbf{Aspecto} & \textbf{Régimen anterior} & \textbf{Régimen actual} \\
\midrule
\textbf{Acceso a la justicia} & Había que pagar la multa para poder impugnarla judicialmente. & Se prohíbe condicionar la revisión judicial al pago previo de la multa. \\
\textbf{Principio de inocencia} & Se presumía la culpabilidad: la pena quedaba aplicada de antemano. & La multa no se cobra hasta que se resuelva la impugnación. \\
\textbf{Alcance} & Muchas normas lo preveían, entre ellas la Ley de Defensa del Consumidor. & Vale para defensa del consumidor, defensa de la competencia y lealtad comercial. \\
\bottomrule
\end{tabularx}
\end{table}

\begin{important}
Exigir el pago previo de la multa para recurrirla afecta el principio de inocencia y el acceso a la justicia; hoy la Ley de Procedimientos Administrativos lo prohíbe de forma genérica.
\end{important}

\section{Lealtad comercial}

\subsection{Concepto y diferencia con la defensa del consumidor}

\begin{definition}{Lealtad comercial}
Subgrupo que sanciona determinados incumplimientos de empresas con giro comercial que venden productos y servicios en el país. El foco no está puesto estrictamente en el consumidor, aunque es el sujeto que más se intenta tutelar. Se sanciona a empresas con información errada o inexacta sobre sus bienes y servicios que les dé algún tipo de ventaja frente a otros competidores.
\end{definition}

Todo ello redunda también en una protección del consumidor, pero el criterio es algo más amplio. Se parece al de la ley de defensa de la competencia, vista en otra clase, que no solo defiende al consumidor sino una suerte de \textbf{juego limpio} en materia comercial, evitando que alguna empresa se favorezca con algún tipo de promoción o rótulo de sus bienes y servicios.

\subsection{El ejemplo de los fideos ``italianos''}

\begin{example}{El paquete de fideos}
Una empresa de fideos coloca en su paquete una banderita italiana y la leyenda ``Tagliatelle, original de Italia'', cuando en realidad se elaboran en una empresa argentina. La empresa juega con una denominación de origen que intenta favorecerla y la hace ver más atractiva, con un producto diferenciado frente a su competencia, y esto incide directamente en la opción del consumidor. Si esos fideos no son de origen italiano, se estaría violando la ley de lealtad comercial, que establece la obligación de informar fehacientemente el origen del producto. Con esta ley se busca desalentar conductas que generen ventajas inexactas frente al consumidor, como invocar el origen italiano, que es un símbolo de estatus y prestigio.
\end{example}

El docente señala que la lealtad comercial era una ley que luego se modificó por decreto, y expresa su opinión de que habría correspondido modificarla también por medio de una ley.

\begin{table}[H]
\centering
\caption{Defensa del consumidor, lealtad comercial y defensa de la competencia}
\small
\begin{tabularx}{\textwidth}{
    >{\raggedright\arraybackslash}p{0.2\textwidth}
    >{\raggedright\arraybackslash}X
    >{\raggedright\arraybackslash}X
}
\toprule
\textbf{Régimen} & \textbf{Foco principal} & \textbf{Ejemplo del docente} \\
\midrule
\textbf{Defensa del consumidor} & Tutela estrictamente al consumidor. & Oferta masiva de carne sin informar el stock. \\
\textbf{Lealtad comercial} & Sanciona a empresas con información errada o inexacta que les da ventaja frente a competidores; protege también al consumidor. & Fideos ``Tagliatelle original de Italia'' que no son de origen italiano. \\
\textbf{Defensa de la competencia} & Defiende el juego limpio en materia comercial, no solo al consumidor. & Evitar que una empresa se favorezca con rótulos o promociones. \\
\bottomrule
\end{tabularx}
\end{table}

\begin{important}
La lealtad comercial obliga a informar fehacientemente, por ejemplo, el origen del producto; protege el juego limpio entre competidores además de proteger al consumidor.
\end{important}

\section{El caso de la factura eléctrica y el Tribunal Fiscal}

\subsection{Multa frente a error de facturación}

\begin{example}{El medidor mal registrado}
Se planteó el caso de una usuaria de una distribuidora eléctrica (Edenor) cuyo medidor marcaba mal. Se le aplicó el régimen de ``primero pague y después reclame'', y recorrió todo el camino administrativo hasta el ENRE, que finalmente le devolvió lo pagado. En la factura se aplicaba además una multa de subida de categoría. El docente explica que la prohibición del \textit{solve et repete} rige desde hace 2 años, con la modificación.
\end{example}

El docente distingue dos situaciones. La prohibición está destinada a evitar que se deba pagar una \textbf{sanción} que se impone para poder discutirla. Un problema de facturación es otra cosa: no es necesariamente una sanción que obliguen a cumplir, sino un diferendo. No es tradicionalmente una violación de la presunción de inocencia, sino una suerte de presunción de que el cálculo de la factura es correcto hasta que se demuestre lo contrario.

Si la factura informa un consumo que aumenta más de un 75\% sobre el consumo promedio para esa época del año, el usuario tiene derecho a no pagar y cuestionarlo, y solo si el cuestionamiento es ratificado queda obligado a pagar. Respecto del elemento sancionatorio de subir de categoría, que se cobra más, habría que ver el caso y el marco regulatorio del ENRE.

% TODO: contenido incompleto o ambiguo en las notas originales (el docente no resuelve si el caso particular encuadraba en la regla del 75\% ni cómo opera el marco regulatorio del ENRE)

\begin{table}[H]
\centering
\caption{Sanción impuesta frente a error de facturación}
\small
\begin{tabularx}{\textwidth}{
    >{\raggedright\arraybackslash}p{0.22\textwidth}
    >{\raggedright\arraybackslash}X
    >{\raggedright\arraybackslash}X
}
\toprule
\textbf{Situación} & \textbf{Naturaleza} & \textbf{Tratamiento según el docente} \\
\midrule
\textbf{Multa o sanción de la administración} & Ejercicio del poder punitivo. & No puede exigirse su pago como condición para recurrir (principio de inocencia). \\
\textbf{Error de facturación} & Diferendo sobre el cálculo; presunción de que la factura es correcta hasta que se demuestre lo contrario. & No es tradicionalmente una violación de la presunción de inocencia; si el consumo aumenta más de un 75\%, se puede no pagar y pedir revisión. \\
\textbf{Multas tributarias} & Sanciones impositivas apelables ante el Tribunal Fiscal. & No se cobran hasta que el tribunal se expida: efecto suspensivo. \\
\bottomrule
\end{tabularx}
\end{table}

\subsection{Un paralelo con el derecho tributario}

En materia fiscal se pueden apelar las determinaciones de oficio que hace ARCA ante el \textbf{Tribunal Fiscal}, un tribunal administrativo que se desempeña en el marco del Ministerio de Economía. Se puede cuestionar una determinación de oficio, es decir, una determinación impositiva, pero todo lo que sean sanciones no se cobra hasta que el tribunal se expida: tiene \textbf{efecto suspensivo} para las multas tributarias. El docente relaciona esto con el mismo concepto del \textit{solve et repete}.

Según la ley actual de procedimientos administrativos, el \textit{solve et repete} ya no existe, y no sería legítima la pretensión de cobrar multas para acceder a la justicia. Además, están en juego la presunción de inocencia y el derecho de acceso a la justicia: el ciudadano no debe encontrar impedimentos ni obstáculos para acceder a la justicia impugnando esa sanción, es decir, debe poder acceder a la vía judicial sin trabas.

\begin{important}
En materia sancionatoria ya no es legítimo exigir el pago de la multa para acceder a la justicia: están en juego la presunción de inocencia y el acceso a la vía judicial sin obstáculos. Un diferendo de facturación es distinto de una sanción.
\end{important}

\section{Cuadro Cornell del capítulo}

\begin{longtable}{
    >{\raggedright\arraybackslash}p{0.27\textwidth}
    !{\vrule width 0.4pt}
    >{\raggedright\arraybackslash}p{\dimexpr0.73\textwidth-6\tabcolsep-0.4pt\relax}
}
\toprule
\textbf{Preguntas / palabras clave} & \textbf{Notas / respuestas} \\
\midrule
\endfirsthead
\toprule
\textbf{Preguntas / palabras clave} & \textbf{Notas / respuestas} \\
\midrule
\endhead
\bottomrule
\endfoot
\bottomrule
\endlastfoot

\textbf{¿Se puede obligar a pagar una multa para poder impugnarla ante la justicia (\textit{solve et repete})?} &
La reforma de 2014 lo preveía. Afecta el principio de inocencia (la pena queda aplicada de antemano) y el acceso a la justicia, además de agredir el patrimonio. La Ley de Procedimientos Administrativos prohíbe de forma genérica estos dispositivos (autorrestricción del Estado). Vale también para defensa de la competencia y lealtad comercial. \\[0.6em]

\textbf{¿Qué protege la lealtad comercial además del consumidor?} &
Sanciona la información errada o inexacta que da ventaja frente a competidores y protege el juego limpio en materia comercial. Obliga a informar fehacientemente el origen del producto. Ejemplo: fideos ``Tagliatelle, original de Italia'' que no son de origen italiano. \\[0.6em]

\textbf{¿Es lo mismo pagar una multa para apelar que pagar una factura que se discute?} &
No. La multa es una sanción y su pago previo no puede exigirse. Un error de facturación es un diferendo, con presunción de que el cálculo es correcto hasta prueba en contrario; si el consumo aumenta más de 75\%, se puede no pagar y pedir revisión. \\[0.6em]

\textbf{¿Qué efecto tiene apelar ante el Tribunal Fiscal una sanción tributaria?} &
Las multas tributarias apeladas ante el Tribunal Fiscal no se cobran hasta que el tribunal se expida: tiene efecto suspensivo. El docente lo relaciona con el mismo concepto del \textit{solve et repete}. \\

\midrule
\multicolumn{2}{p{\dimexpr\textwidth-4\tabcolsep\relax}}{%
\textbf{Resumen:} exigir el pago previo de una multa para impugnarla (\textit{solve et repete}) afecta el principio de inocencia y el acceso a la justicia, y hoy está prohibido de forma genérica por la Ley de Procedimientos Administrativos. La lealtad comercial protege el juego limpio entre competidores y obliga a informar fehacientemente datos como el origen del producto. Una sanción no es lo mismo que un diferendo de facturación, y las multas tributarias apeladas ante el Tribunal Fiscal no se cobran hasta que este se expida.
} \\
\end{longtable}

\end{document}
